% GENERATED FILE. Edit canonical lessons, then run npm run generate:books.
% canonical-source-hash: fnv1a64:4dc56b8b563e0585
% canonical-lessons: JA-C433-akogareru, JA-C433-koishii, JA-C433-urayamashii, JA-C433-kyoumi, JA-C433-kanshin

\chapter{Longing and Envy}
\label{ch:ja-longing-and-envy}
\hlchaptermodality{433}

\begin{chapteropening}
\textbf{By the end of this chapter:} \emph{I can say what I long for and what interests me.}

Complete the last lesson of chapter 433: I can say what I long for and what interests me.
\end{chapteropening}

\section[\texorpdfstring{\ja{あこがれる} (akogareru)}{あこがれる (akogareru)}]{\ja{あこがれる} (akogareru) — to yearn after; to be attracted by}

\label{lesson:JA-C433-akogareru}



\begin{quote}
Before the new one: say the Japanese for futile, then the Japanese for to give up.
\end{quote}

\subsection*{You'll want to know: \ja{あこがれる}}
\textbf{\ja{あこがれる}} — \emph{akogareru} — \textquotedblleft{}to yearn after; to be attracted by\textquotedblright{}.

\begin{quote}
\textbf{\ja{おとうとは} \ja{せんしゅに} \ja{あこがれています。}} — \emph{otōto wa senshu ni akogarete imasu} — \textquotedblleft{}My younger brother admires the athletes.\textquotedblright{}
\end{quote}

Polite \textbf{\ja{あこがれます}}, \emph{akogaremasu}; an ichidan verb. Past (polite): \textbf{\ja{むすこが} \ja{せんぱいに} \ja{あこがれました。}}, \emph{musuko ga senpai ni akogaremashita}.

\begin{cousinweb}
\textbf{\ja{あこがれる}} (\emph{akogareru}) is a native Japanese word (wago), inherited rather than borrowed from Chinese; its kanji spelling was attached to it later (JMdict).
\end{cousinweb}

\begin{grammarlens}[title={What you've built}]
One more word for describing what you lived and felt.
\end{grammarlens}

\subsection*{Your turn}
\begin{itemize}
\raggedright
  \item \emph{Say it:} \emph{akogareru}
  \item \emph{Say it:} \emph{akogareru}, once more
  \item \emph{Say it:} \emph{otōto wa senshu ni akogarete imasu} — My younger brother admires the athletes.
  \item \emph{Recall:} say the Japanese for futile, then the Japanese for to give up, then say \emph{akogareru} again
\end{itemize}

\begin{tcolorbox}[breakable,colback=teal!4,colframe=teal!35!black,title={Before you move on}]
What does \ja{あこがれる} mean? (\textquotedblleft{}To yearn after; to be attracted by\textquotedblright{}.) Say it once more.
\end{tcolorbox}
\section[\texorpdfstring{\ja{こいしい} (koishii)}{こいしい (koishii)}]{\ja{こいしい} (koishii) — longed for; missed dearly}

\label{lesson:JA-C433-koishii}



\begin{quote}
Before the new one: say the Japanese for to give up, then the Japanese for to yearn after.
\end{quote}

\subsection*{You'll want to know: \ja{こいしい}}
\textbf{\ja{こいしい}} — \emph{koishii} — \textquotedblleft{}longed for; missed dearly\textquotedblright{}.

\begin{quote}
\textbf{\ja{ふゆは} \ja{あたたかい} \ja{なつが} \ja{こいしいです。}} — \emph{fuyu wa atatakai natsu ga koishii desu} — \textquotedblleft{}In winter I long for the warm summer.\textquotedblright{}
\end{quote}

An i-adjective.

\begin{cousinweb}
\textbf{\ja{こいしい}} (\emph{koishii}) is a native Japanese word (wago), inherited rather than borrowed from Chinese; its kanji spelling was attached to it later (JMdict).
\end{cousinweb}

\begin{grammarlens}[title={What you've built}]
One more word for describing what you lived and felt.
\end{grammarlens}

\subsection*{Your turn}
\begin{itemize}
\raggedright
  \item \emph{Say it:} \emph{koishii}
  \item \emph{Say it:} \emph{koishii}, once more
  \item \emph{Say it:} \emph{fuyu wa atatakai natsu ga koishii desu} — In winter I long for the warm summer.
  \item \emph{Recall:} say the Japanese for to give up, then the Japanese for to yearn after, then say \emph{koishii} again
\end{itemize}

\begin{tcolorbox}[breakable,colback=teal!4,colframe=teal!35!black,title={Before you move on}]
What does \ja{こいしい} mean? (\textquotedblleft{}Longed for; missed dearly\textquotedblright{}.) Say it once more.
\end{tcolorbox}
\section[\texorpdfstring{\ja{うらやましい} (urayamashii)}{うらやましい (urayamashii)}]{\ja{うらやましい} (urayamashii) — envious}

\label{lesson:JA-C433-urayamashii}



\begin{quote}
Before the new one: say the Japanese for to yearn after, then the Japanese for missed dearly.
\end{quote}

\subsection*{You'll want to know: \ja{うらやましい}}
\textbf{\ja{うらやましい}} — \emph{urayamashii} — \textquotedblleft{}envious\textquotedblright{}.

\begin{quote}
\textbf{\ja{あねの} \ja{くるまが} \ja{うらやましいです。}} — \emph{ane no kuruma ga urayamashii desu} — \textquotedblleft{}I envy my elder sister her car.\textquotedblright{}
\end{quote}

An i-adjective.

\begin{cousinweb}
\textbf{\ja{うらやましい}} (\emph{urayamashii}) is a native Japanese word (wago), inherited rather than borrowed from Chinese; its kanji spelling was attached to it later (JMdict).
\end{cousinweb}

\begin{grammarlens}[title={What you've built}]
One more word for describing what you lived and felt.
\end{grammarlens}

\subsection*{Your turn}
\begin{itemize}
\raggedright
  \item \emph{Say it:} \emph{urayamashii}
  \item \emph{Say it:} \emph{urayamashii}, once more
  \item \emph{Say it:} \emph{ane no kuruma ga urayamashii desu} — I envy my elder sister her car.
  \item \emph{Recall:} say the Japanese for to yearn after, then the Japanese for missed dearly, then say \emph{urayamashii} again
\end{itemize}

\begin{tcolorbox}[breakable,colback=teal!4,colframe=teal!35!black,title={Before you move on}]
What does \ja{うらやましい} mean? (\textquotedblleft{}Envious\textquotedblright{}.) Say it once more.
\end{tcolorbox}
\section[\texorpdfstring{\ja{きょうみ} (kyōmi)}{きょうみ (kyōmi)}]{\ja{きょうみ} (kyōmi) — an interest; curiosity}

\label{lesson:JA-C433-kyoumi}



\begin{quote}
Before the new one: say the Japanese for missed dearly, then the Japanese for envious.
\end{quote}

\subsection*{You'll want to know: \ja{きょうみ}}
\textbf{\ja{きょうみ}} — \emph{kyōmi} — \textquotedblleft{}an interest; curiosity\textquotedblright{}.

\begin{quote}
\textbf{\ja{おんがくに} \ja{きょうみが} \ja{あります。}} — \emph{ongaku ni kyōmi ga arimasu} — \textquotedblleft{}I am interested in music.\textquotedblright{}
\end{quote}

\begin{cousinweb}
\textbf{\ja{きょうみ}} (\emph{kyōmi}) is Sino-Japanese (kango): two Chinese characters read with their on'yomi, readings borrowed from Chinese. They are \emph{kyō}, \textquotedblleft{}interest\textquotedblright{} (Middle Chinese \emph{xingH}) and \emph{mi}, \textquotedblleft{}taste\textquotedblright{} (\emph{mj+jH}) (KANJIDIC2; Middle Chinese in Baxter's transcription).
\end{cousinweb}

\begin{grammarlens}[title={What you've built}]
One more word for describing what you lived and felt.
\end{grammarlens}

\subsection*{Your turn}
\begin{itemize}
\raggedright
  \item \emph{Say it:} \emph{kyōmi}
  \item \emph{Say it:} \emph{kyōmi}, once more
  \item \emph{Say it:} \emph{ongaku ni kyōmi ga arimasu} — I am interested in music.
  \item \emph{Recall:} say the Japanese for missed dearly, then the Japanese for envious, then say \emph{kyōmi} again
\end{itemize}

\begin{tcolorbox}[breakable,colback=teal!4,colframe=teal!35!black,title={Before you move on}]
What does \ja{きょうみ} mean? (\textquotedblleft{}An interest; curiosity\textquotedblright{}.) Say it once more.
\end{tcolorbox}
\section[\texorpdfstring{\ja{かんしん} (kanshin)}{かんしん (kanshin)}]{\ja{かんしん} (kanshin) — admiration; being impressed}

\label{lesson:JA-C433-kanshin}



\begin{quote}
Before the new one: say the Japanese for envious, then the Japanese for an interest.
\end{quote}

\subsection*{You'll want to know: \ja{かんしん}}
\textbf{\ja{かんしん}} — \emph{kanshin} — \textquotedblleft{}admiration; being impressed\textquotedblright{}.

\begin{quote}
\textbf{\ja{かれの} \ja{しごとに} \ja{かんしんしました。}} — \emph{kare no shigoto ni kanshin shimashita} — \textquotedblleft{}I was impressed by his work.\textquotedblright{}
\end{quote}

\begin{cousinweb}
\textbf{\ja{かんしん}} (\emph{kanshin}) is Sino-Japanese (kango): two Chinese characters read with their on'yomi, readings borrowed from Chinese. They are \emph{kan}, \textquotedblleft{}emotion\textquotedblright{} (Middle Chinese \emph{komX}) and \emph{shin}, \textquotedblleft{}heart\textquotedblright{} (\emph{sim}) (KANJIDIC2; Middle Chinese in Baxter's transcription).
\end{cousinweb}

\begin{grammarlens}[title={What you've built}]
One more word for describing what you lived and felt.
\end{grammarlens}

\subsection*{Your turn}
\begin{itemize}
\raggedright
  \item \emph{Say it:} \emph{kanshin}
  \item \emph{Say it:} \emph{kanshin}, once more
  \item \emph{Say it:} \emph{kare no shigoto ni kanshin shimashita} — I was impressed by his work.
  \item \emph{Recall:} say the Japanese for envious, then the Japanese for an interest, then say \emph{kanshin} again
\end{itemize}

\begin{tcolorbox}[breakable,colback=teal!4,colframe=teal!35!black,title={Before you move on}]
What does \ja{かんしん} mean? (\textquotedblleft{}Admiration; being impressed\textquotedblright{}.) Say it once more.
\end{tcolorbox}
